\section{Related Work}
\label{sec:related_work}

A rigorous body of literature underpins the rapid evolution of Earth Observation (EO) foundation models. We categorize these advancements into agricultural monitoring applications, masked image modeling paradigms, and the emergent shift toward hybrid contrastive-reconstructive frameworks.

\subsection{Agricultural and Aerial Infrastructure Detection}
The automated mapping of agricultural infrastructure---such as concentrated animal feeding operations (CAFOs) and greenhouses---remains a critical environmental and regulatory target. Historically, these efforts have relied heavily on supervised Convolutional Neural Networks (CNNs) initialized with ImageNet priors \cite{wang2022empirical, robinson2021mapping}. For instance, supervised object detectors like YOLO and Faster R-CNN have been deployed extensively for greenhouse detection using high-resolution multispectral imagery \cite{li2020greenhouse}. Similarly, models have been scaled to process terabytes of National Agriculture Imagery Program (NAIP) data for continental CAFO mapping \cite{robinson2021mapping}, and have seen broad adoption in ecological benchmarks like DOTA \cite{xia2017dota, ding2021object} and tree-crown health detection \cite{joshi2025vision}. However, these approaches suffer from poor generalization when deployed across heterogeneous topographies, sensors, and temporal gaps due to the profound domain mismatch between ground-level natural imagery and top-down aerial photography \cite{wang2022empirical, marsocci2024cross}.

\subsection{Masked Image Modeling in Remote Sensing}
To circumvent the bottleneck of dense annotations, the EO community has increasingly adopted Self-Supervised Learning (SSL). Masked Image Modeling (MIM), inspired by BEiT \cite{bao2021beit} and instantiated in architectures like SparK \cite{tian2023spark}, forces the model to reconstruct occluded patches. This paradigm excels at capturing local spatial dependencies and has become the dominant foundation-model paradigm for multimodal EO datasets \cite{sun2023ringmo, wang2024scaling}. For example, RingMo \cite{sun2023ringmo} and Scale-MAE \cite{reed2022scale} successfully scaled MIM to millions of remote sensing images, achieving state-of-the-art results across various dense prediction tasks. Similarly, cross-scale and multimodal MAE variants \cite{tang2024cross, sosa2025multimae} have demonstrated strong parameter efficiency and transferability.

Nevertheless, vanilla MIM is not universally optimal. A growing consensus argues that random high-ratio masking often prioritizes pixel-level local fidelity at the expense of global semantic coherence. In dense, cluttered remote sensing scenes, this can erase small objects and fine local structures, leading to suboptimal embeddings for linearly separable tasks or zero-shot inference \cite{nedungadi2024mmearth}.

\subsection{Contrastive Self-Distillation and Hybrid Models}
In contrast to purely reconstructive objectives, contrastive learning frameworks focus on aligning representations across augmented views, ensuring a distinct and linearly-separable feature space \cite{jain2022multimodal}. The DINO (self-distillation with no labels) framework \cite{caron2021emerging} optimizes a cross-entropy objective between teacher and student networks to learn powerful global semantics. Several studies demonstrate that contrastive alignment yields superior cross-sensor semantic alignment compared to vanilla MAE \cite{wolf2026brewing, guo2023skysense}. For instance, SkySense scaled multimodal contrastive pretraining to a billion-parameter level, demonstrating that domain-specific foundation models can surpass generic ImageNet variants across numerous benchmarks \cite{guo2023skysense}.

However, because pure contrastive objectives are susceptible to dimensional collapse, the strongest recent EO systems increasingly deploy hybrid designs. Innovations such as KoLeo regularization \cite{sablayrolles2021spreading} and hybridized objectives like iBOT \cite{zhou2021ibot} integrate masked patch prediction with self-distillation. Similarly, models like CROMA \cite{fuller2023croma}, CMID \cite{muhtar2023cmid}, SegMind \cite{li2023segmind}, and MaCon \cite{wang2025macon} explicitly demonstrate that contrastive and reconstructive objectives are highly complementary. This hybridization recovers fine-grained local structures while strictly regularizing the global embedding manifold, setting a new standard for robust transfer learning in Earth Observation.
